%%%%%%%%%%%%%%%%%%%%%%%%%%%%%%%%%%%%%%%%%%%%%%%%%%%%%%%%%%%%%%%%%%%%%%%%%%%
% Handout und Uebungsblatt zu Lehreinheit 6
% Im Skript am Kapitelende, in handouts.tex als Druckvorlage.
%%%%%%%%%%%%%%%%%%%%%%%%%%%%%%%%%%%%%%%%%%%%%%%%%%%%%%%%%%%%%%%%%%%%%%%%%%%
\begin{handout}{Forschungsdesign}

\textbf{Worum es heute ging.} Drei Seile wurden gleich lang, zweimal. Beim
ersten Mal gab es nur die Vorführung: Wer behauptete, zeigte auch, und
niemand maß. Beim zweiten Mal legten Sie vorher fest, wie geprüft wird.
Ohne Vergleich bleiben selbst grobe Täuschungen unentdeckt; eine
verweigerte Kontrolle ist selbst ein Befund.

\textbf{Kernbegriffe}
\begin{description}
  \item[Forschungsdesign] Plan einer Untersuchung: welche Variablen an
    wem, wann und unter welchen Bedingungen erhoben oder verändert werden
    und womit verglichen wird.
  \item[UV / AV / Störvariable] vermutete Ursache / gemessene Wirkung /
    weiterer Einfluss auf die AV.
  \item[Konfundierung] eine Störvariable variiert systematisch mit der
    UV; ihre Wirkung ist nicht mehr abtrennbar.
  \item[Interne Validität] die Veränderung der AV ist eindeutig auf die
    UV zurückzuführen \parencite{shadish2002experimental}.
  \item[Externe Validität] das Ergebnis lässt sich auf andere Personen,
    Situationen und Zeiten übertragen.
  \item[Objektivität, Reliabilität, Validität] unabhängig von der
    messenden Person; bei Wiederholung gleich; misst das Gemeinte.
  \item[Blind-Design, Parallelmessung] Messende kennen die Bedingung
    nicht; dasselbe Konstrukt wird mit mehreren Verfahren erhoben.
\end{description}

\textbf{Designtypen}

\begin{center}
\small
\begin{tabular}{p{0.24\textwidth}p{0.14\textwidth}p{0.1\textwidth}p{0.36\textwidth}}
\toprule
& \textbf{UV variiert} & \textbf{Zufall} & \textbf{Kurzschrift (R, X, O)} \\
\midrule
Experiment & ja & ja & R\,X\,O \;/\; R\;\;\,O \\
Quasi-Experiment & ja & nein & \phantom{R\,}X\,O \;/\; \phantom{R\,X\,}O \\
nicht-experimentell & nein & -- & nur O (Zusammenhänge) \\
\bottomrule
\end{tabular}
\end{center}

\textbf{Bedrohungen der internen Validität:} Geschichte, Reifung,
Testeffekt, Instrumentierung, Regression zur Mitte, Selektion, Ausfall
-- und die Erwartungen der Forschenden.

\textbf{Fünf Werkzeuge der Kontrolle:} Randomisierung, Konstanthalten,
Schichten/Parallelisieren, (aktive) Kontrollgruppe, Blindbedingungen.

\begin{merkbox}
Ein Design ist so gut wie die Alternativerklärungen, die es ausschließt.
Fragen Sie bei jeder Kontrollgruppe: Unterscheidet sie sich \emph{nur} in
der UV?
\end{merkbox}

\textbf{Zum Weiterlesen:} \textcite{doering2016forschungsmethoden},
Kapitel zu Untersuchungsdesign und Gütekriterien;
\textcite{shadish2002experimental}, Kap.~1--2.
\end{handout}

\begin{uebungsblatt}{Forschungsdesign}

\aufgabe{Fallstudie: Ein Design analysieren}[in der Stunde, Zweiergruppe, 15 Min.]
\textit{Die folgende Studie ist fiktiv und für diese Übung konstruiert;
alle Zahlen sind erfunden.}

\begin{center}
\fbox{\parbox{0.92\textwidth}{\small
\textbf{\enquote{Lern-App verbessert Mathematiknoten um 0,7}}\\[1mm]
Im Wintersemester bietet eine Hochschule zur Vorlesung Mathematik~1 eine
neue Übungs-App an, die der Dozent der Vorlesung selbst entwickelt hat.
Die Nutzung ist freiwillig; 64 von 212 Studierenden melden sich an. Zur
Einführung der App bietet der Dozent für die Angemeldeten ein
wöchentliches Tutorium an, in dem mit der App geübt wird. Am
Semesterende schreiben alle 212 Studierenden dieselbe Klausur, die der
Dozent selbst korrigiert; auf den Bögen stehen Name und Matrikelnummer.
Die App-Gruppe erreicht im Mittel die Note 2,3, die übrigen 148
Studierenden die Note 3,0. Im Pressetext der Hochschule heißt es:
\enquote{Die App verbessert die Note um 0,7. Sie wird ab dem nächsten
Semester für alle Erstsemester verpflichtend.}
}}
\end{center}

\begin{enumerate}
  \item Benennen Sie UV und AV (unabhängige und abhängige Variable). Wie
    ist die AV gemessen worden?
  \item Welcher Designtyp liegt vor? Schreiben Sie das Design in der
    Kurzschrift R/X/O (eine Zeile je Gruppe; R = zufällige Zuteilung,
    X = Behandlung, O = Messung).
  \item Gibt es eine Kontrollgruppe? In welchen Merkmalen unterscheidet
    sie sich von der App-Gruppe außer in der App-Nutzung?
  \item Nennen Sie mindestens drei Konfounder bzw.\ Bedrohungen der
    internen Validität und ordnen Sie sie Fachbegriffen zu (etwa Selektion,
    Geschichte, Reifung, Testeffekt, Instrumentierung, Ausfall,
    Erwartungen der Forschenden).
  \item Ist die Messung objektiv? Was wäre ein Blind-Design?
  \item Skizzieren Sie ein verbessertes Design, das die Aussage des
    Pressetexts prüfen könnte.
\end{enumerate}
\schreiblinien{9}

\aufgabe{Design-Skizze zu Ihrer Hypothese}[Beginn in der Stunde, Abgabe bis LE~7, benotet]
Füllen Sie die Vorlage aus und begründen Sie Ihre Wahl auf höchstens einer
weiteren Seite.

\begin{center}
\renewcommand{\arraystretch}{1.5}
\begin{tabular}{|p{0.3\textwidth}|p{0.6\textwidth}|}
\hline
H\textsubscript{1} (aus LE~5) & \\[6mm] \hline
UV und ihre Stufen & \\ \hline
AV und Messinstrument & \\ \hline
Designtyp (experimentell / quasi / nicht-experimentell; Quer- oder Längsschnitt) & \\ \hline
Design in Kurzschrift (R / X / O) & \\ \hline
Stichprobe: wer, wie viele, wie rekrutiert? & \\ \hline
Zuteilung zu den Bedingungen & \\ \hline
Kontrollbedingung (aktiv?) & \\ \hline
Wichtigste Störvariablen und wie ich sie kontrolliere & \\[8mm] \hline
Messgüte: objektiv? reliabel? valide? blind? Parallelmessung? & \\[8mm] \hline
Verbleibende Bedrohung, die ich in Kauf nehme & \\ \hline
\end{tabular}
\end{center}
\end{uebungsblatt}

\begin{uebungsblattb}{3}{Feedback und Transfer}

\aufgabe{Peer-Feedback}[in der Stunde, Zweiergruppe, 10 Min.]
Notieren Sie für die Skizze Ihrer Partnerin bzw.\ Ihres Partners:
\begin{center}
\renewcommand{\arraystretch}{2}
\begin{tabular}{|p{0.42\textwidth}|p{0.48\textwidth}|}
\hline
Eine Alternativerklärung, die das Design nicht ausschließt & \\ \hline
Eine Änderung, die sie ausschließen würde & \\ \hline
Was an dem Design überzeugt & \\ \hline
\end{tabular}
\end{center}

\aufgabe{Tagebuch des Staunens}[Hausaufgabe]
Finden Sie in Werbung, Nachrichten oder sozialen Medien eine
Wirkungsbehauptung, der ein Vergleich fehlt. Welche Kontrollbedingung
bräuchte man, um sie zu prüfen?

\begin{lehrhinweis}[Erwartungshorizont]
\textbf{Aufgabe~1.} UV: App-Nutzung (ja/nein, selbst gewählt); AV:
Klausurnote, korrigiert vom Entwickler der App, nicht anonymisiert. Design:
quasi-experimentell bzw.\ nicht-experimenteller Gruppenvergleich ohne
Vortest (X\,O / \phantom{X\,}O), keine Randomisierung. Die
\enquote{Kontrollgruppe} unterscheidet sich in Motivation und vermutlich
Vorkenntnissen (Selbstselektion), im zusätzlichen Tutorium (Konfundierung
der UV mit Übungszeit und Betreuung) und darin, wie der Korrektor sie
wahrnimmt (Versuchsleitereffekt, fehlende Objektivität). Weitere Punkte:
kein Vortest, also keine Aussage über Veränderung; mögliche
Neuheitseffekte; Übertragung von Freiwilligen auf eine Pflichtnutzung ist
fraglich (externe Validität). Verbessertes Design: zufällige Zuteilung
(ggf.\ geschichtet nach Vorkenntnissen aus einem Eingangstest), Tutorium
für beide Gruppen gleich lang mit herkömmlichen Übungen (aktive
Kontrolle), anonymisierte Klausur, Korrektur durch eine Person, die die
Gruppenzugehörigkeit nicht kennt (Blind-Design).
\textbf{Aufgabe~2.} Bewertet werden: Passung von Design und Hypothese
(Kausalbehauptung erfordert experimentelles oder gut begründetes
quasi-experimentelles Design), eine Kontrollbedingung, die sich nur in der
UV unterscheidet, ein konkretes Messinstrument, mindestens zwei benannte
und kontrollierte Störvariablen, eine ehrlich benannte Restbedrohung.
\end{lehrhinweis}
\end{uebungsblattb}
